\section{Follow-up 1: the four rows in Grothendieck's frameworks}
\label{sec:f1-grothendieck}

This section records the realization maps for HD-F1-T01--T09,
HD-F1-S01--S11 and HD-F1-R01--R11, followed by the explicitly
dated extensions for subsequent findings. Its assertions are consequences
of the constructions proved in those sections. The source locators
below identify the formalism; they are not claims that the particular
degree computations occur in those sources. The four DEG\_L row names
are retained throughout.

\subsection{HD-F1-F01: flat classes, their actual cycle families,
and the variational statement}

Let $\pi:\mathcal A\to S$ be one of the smooth polarized Weil
families actually used above, after a finite level cover when needed
to carry its endomorphisms and polarization line bundle. On a
simply connected analytic chart, the algebraic de Rham bundle
\[
 \mathcal H=R^{2n}\pi_*\Omega^\bullet_{\mathcal A/S},
 \qquad \nabla:\mathcal H\longrightarrow
                       \mathcal H\otimes\Omega_S^1
 \tag{F1F1.1}
\]
has its Betti comparison. Write
$\mathcal W_B=\bigwedge_K^{2n}R^1\pi_*\mathbb Q$
for the Weil local system and
$\mathcal W_{\mathrm{dR}}=\mathcal W_B\otimes\mathcal O_S$
for its flat subbundle under that comparison.
The actual graph correspondences defining $e_W$ commute with
Gauss--Manin transport: a relative endomorphism acts on the relative
de Rham complex, and its comparison with the Betti local system
is natural. In particular $e_W$ is horizontal.

For any flat family of the integral cycles constructed here,
its cycle class is horizontal and lies in $F^n\mathcal H$.
The cycle-class compatibility is witnessed by that family itself.
For a proper isogeny $h$ and an integral source component $Z$
with reduced image $Y$, the exact comparison in every fibre is
\[
 \operatorname{cl}(h_*Z)
       =h_*\operatorname{cl}(Z)
       =\deg(Z/Y)\operatorname{cl}(Y),\qquad
 e_Wh_*=h_*e_W.
 \tag{F1F1.2}
\]
On the three-dimensional graph Hodge locus computed in
HD-F1-S03, the horizontal section supplied by the proper push in
HD-F1-T02 is $p^{2k}[Y_k]$. Its reduced image supplies the actual
integral class $[Y_k]$ of degree eight, with the displayed integer
relating the two. In the Schoen construction of HD-F1-S07 the
relative degree-nine quotient supplies the degree-at-most-36
detecting image on a dominant family. The already checked proper
cycle-space specialization extends that bound throughout its
principal component. These are the actual variational realizations
for DEG\_L Q-BOUNDS and Q-DIV.

Grothendieck's original formulation is in \emph{On the de Rham
cohomology of algebraic varieties}, Publ. Math. IHES 29 (1966),
p.103, footnote13, added in July1965; the de Rham comparison is
Theorem1$'$ on p.96. The formulation uses a horizontal relative
de Rham class and its algebraicity on a fibre. Here the particular
cycle families and the bounded-cycle specialization have been
constructed and checked. No general variational conjecture is
being added as a premise.

For DEG\_L Q-T7-FIXED, HD-F1-S03 computes the full class of the
graph. Its Hodge locus is the image of the explicit morphism
\[
 \operatorname{Sp}(U,\psi_U)\longrightarrow U(2,2),
                     \qquad g\longmapsto\operatorname{diag}(g,g).
 \tag{F1F1.3}
\]
It has dimension three. In the fourth tangent direction of the
Weil domain, the obstruction is the nonzero skew part
$i(Z^t-Z)/(2\sqrt3)$, even though $e_W[Y]$ is horizontal of
type$(2,2)$. Formula(F1F1.3), the contraction-kernel recovery of
$U$ from $[Y]$, and the actual obstruction map give the precise
relation between these two class conditions.

\subsection{HD-F1-F02: Hilbert points, Chow images and semiregularity}

For each graph in HD-F1-T02, restriction of the actual polarization
line bundle gives $L|_Y$ with $c_1(L|_Y)=2\beta$ and
$(L|_Y)^2=8$. Riemann--Roch on its abelian surface therefore gives
the fixed Hilbert polynomial
\[
 \chi(Y,L^{\otimes r}|_Y)=4r^2.
 \tag{F1F2.1}
\]
Its translations give a morphism $A/Y\to\operatorname{Hilb}_{4r^2}(A)$,
with tangent map the identity
$\operatorname{Lie}(A/Y)\simeq H^0(Y,N_{Y/A})$.
Over the three-dimensional domain in(F1F1.3), the fixed rational
subspace $U$ and its intersection with the marked lattice give
the corresponding family of abelian subvarieties and Hilbert points.
On the fixed source graph $\Gamma_0$ in HD-F1-T02, the pullback
polarization has class $p^k\theta_0|_{\Gamma_0}$, hence its Hilbert
polynomial is $4p^{2k}r^2$. Proper pushforward on cycle spaces
then sends this source to $p^{2k}[Y_k]$, with its actual
multiplicity, as in(F1F1.2).

For a smooth surface $i:Y\hookrightarrow A$, the tangent-normal
sequence and adjunction supply the exact dual maps
\[
 H^1(A,T_A)\xrightarrow{\operatorname{ob}_Y}H^1(Y,N_{Y/A}),
 \qquad
 H^1(Y,N_{Y/A})\xrightarrow{\sigma_Y}H^3(A,\Omega_A^1),
 \qquad
 \sigma_Y\operatorname{ob}_Y(\kappa)=\kappa\lrcorner[Y].
 \tag{F1F2.2}
\]
The first map is the S-02 conormal map under the fourfold volume
and rank-two determinant identifications of HD-F1-S01.
HD-F1-S02 proves its surjectivity for the graphs, and HD-F1-S05
does so for the specified open smooth Schoen members. For the
actual connected multiplication pullbacks $S_m$, HD-F1-S06 gives
\[
 \dim\operatorname{coker}\mu_{S_m}\ge4m^8-8>0,
 \qquad [S_m]=m^4[S],\qquad \deg_\beta S_m=36m^4.
 \tag{F1F2.3}
\]
Pullback of the smooth Schoen family is an actual Hilbert family
of these surfaces. Thus the failure in(F1F2.3) is compatible with
their deformation, through an explicit cartesian square; the
calculation makes no converse inference from deformation to
injectivity of the semiregularity map.

The exact projective parameter-space statements used here are
Grothendieck, FGA, Bourbaki Expose221, Theorems3.1--3.2,
pp.259--260, for fixed-Hilbert-polynomial Quot schemes and their
Hilbert specialization after Proposition3.8. Proper pushforward
in the displayed cycle calculation is the usual morphism on
effective cycle spaces; the bounded Chow argument remains
ATTACK LemmaDom and HECKE\_DIV R21--R23 at their checked scope.
FGA221 is not being cited for a uniform bound on Hilbert
polynomials merely from a numerical degree: the graph polynomial
has been computed, and the general bounded-degree passage uses
the supplied Chow argument.

For DEG\_L Q-T7-REPLACEMENT, the actual perfect-complex moduli
calculation retains the full trace/evaluation morphisms. The
sixfold object HD-N06 has its explicit coherent finite action,
negative-Ext vanishing, and invariant degree-two trace injection.
HD-F1-R02 gives the analogous small-dimensional control through
its written Fourier kernel. Their fixed whole characters enter
the supplied S-14 deformation and degree-extraction arguments.
The product-orbit obstruction HD-F1-R03 concerns those same full
contraction maps: the product ranks94 and100 are compared through
invertible Hochschild intertwiners with the target ranks
$56,72,88,104$. The additional nonsimple RM restrictions in
HD-F1-R04 concern the actual invariant Ext groups. These object
maps specify what is retained before any passage to a cycle space.

\subsection{HD-F1-F03: the Poincare biextension and degree identities}

Let $\mathcal P_A$ be the normalized Poincare line bundle on
$A\times\widehat A$. For every homomorphism $h:B\to A$, its
dual is defined by pullback of rigidified degree-zero line bundles.
This definition gives the canonical isomorphism of biextensions
on $B\times\widehat A$,
\[
 (h\times1)^*\mathcal P_A
       \simeq(1\times\widehat h)^*\mathcal P_B.
 \tag{F1F3.1}
\]
Indeed both restrictions at a point $\alpha\in\widehat A$ are
the rigidified line bundle $h^*\alpha$ on $B$; the universal
property and rigidification give the isomorphism. For a line
bundle $L$ on $A$, its associated homomorphism consequently obeys
\[
 \lambda_{h^*L}=\widehat h\,\lambda_Lh.
 \tag{F1F3.2}
\]
This is the biextension realization of the actual polarization
pullbacks used above.

For the three maps in the follow-up, these are explicitly
\[
 \begin{split}
 \widehat h_k\lambda_{\theta_k}h_k&=p^k\lambda_{\theta_0},\\
 \widehat j_k\lambda_{\theta_k}j_k&=2\lambda_{\beta_k},\\
 \widehat q\lambda_\theta q&=\lambda_\beta
       \quad(\deg q=9).
 \end{split}
 \tag{F1F3.3}
\]
The last identity is verified by $R^\dagger R=
\left(\begin{smallmatrix}2&-1\\-1&2\end{smallmatrix}\right)$
in HD-F1-S07. The first two give degrees $8p^{2k}$ and8 on
the source and target graph surfaces, respectively. The last
gives degree36 for the proper Schoen image. The identity
$(h_k)_*[\Gamma_0]=p^{2k}[Y_k]$ records the exact divisor by
which the first degree becomes the second.

These statements use SGA7 I, VII, Example2.9.5 and
Remark2.9.6, pp.175--176, especially (2.9.6.1)--(2.9.6.2),
and VIII,3.2, equations(3.2.2)--(3.2.4):
\[
 \operatorname{Biext}^1(P,Q;\mathbb G_m)
       \simeq\operatorname{Hom}(P,\widehat Q).
 \tag{F1F3.4}
\]
The original VII scans were read. The VIII statement was read
in the text transcription and agrees with the original VII
isomorphism just cited; an original VIII scan was not obtained.
The formulas(F1F3.1)--(F1F3.3) were also verified directly above.

\subsection{HD-F1-F04: Adams operations and the retained object maps}

On $K_0(A)$, the Adams operation $\psi^a$, $a\ge1$, is the
Newton polynomial in exterior operations. The splitting principle,
with $\psi^a[L]=[L^{\otimes a}]$ for a line bundle, gives
\[
 \operatorname{ch}_j(\psi^a x)=a^j\operatorname{ch}_j(x),
 \quad
 \operatorname{ch}(\psi^a(e^Bv))
       =e^{aB}\sum_j a^jv_j.
 \tag{F1F4.1}
\]
Here the second equation is notation for the graded character
action; it does not assert that every rational $e^B$ is an
integral $K$-class. The Newton definition and multiplicativity
are SGA6, ExposeV, Appendix7.1--7.5, pp.360--364;
the displayed identities also follow on the flag bundle from
$\sum_i e^{ax_i}$, and injectivity in the splitting principle
brings them back to $A$.

The actual graph class has trivial normal bundle, so GRR gives
$\operatorname{ch}(j_*\mathcal O_Y)=[Y]$. Thus
\[
 \operatorname{ch}(\psi^a[j_*\mathcal O_Y])=a^2[Y],
 \qquad
 \operatorname{ch}((h_k)_*\mathcal O_{\Gamma_0})
                    =p^{2k}[Y_k].
 \tag{F1F4.2}
\]
GRR in the second formula has Todd factors one on both abelian
surfaces and the ambient abelian variety. It is a rational
character identity with the genuine finite pushforward, so it
needs no additional integral Adams/pushforward commutation rule.

For HD-F1-R02, formula(F1F4.1) is the exact identity
\[
 e^{-aP/2}\operatorname{ch}(\psi^a[E])
       =a^2t-2a\theta+\frac{a^3}{12}\theta^3.
 \tag{F1F4.3}
\]
For the carried sixfold object HD-N06 it is
\[
 \operatorname{ch}(\psi^a[\mathcal E_d])
 =32a^3t_d+4-8a^2\theta_0^2
       +\frac83a^4\theta_0^4-\frac{16}{45}a^6\theta_0^6.
 \tag{F1F4.4}
\]
These record the integral operation and its precise graded
realization. The semiregularity statements remain attached to
the actual objects and action maps constructed in HD-N06 and
HD-F1-R02. The forgetful morphism to $K_0$ explains why the
Ext checks have been retained: for every nonzero perfect object
$F$, $[F\oplus F[1]]=0$, whereas
\[
 \operatorname{Ext}^{-1}(F,F[1])
          =\operatorname{Hom}(F,F)\ne0.
 \tag{F1F4.5}
\]
Adding this actual summand preserves a character and destroys
negative-Ext vanishing. This gives a concrete kernel of the
forgetful map and its effect on the object hypothesis.

Finally the displayed morphisms, kernels, finite groups and
cycle families have finite presentation. EGA IV$_3$,8.8.2,
spreads any specified finite collection over a finitely generated
$\mathbb Z$-algebra; shrinking makes the relevant abelian schemes
smooth and the stated finite isogenies etale away from their
degrees. Their Betti, de Rham and $\ell$-adic cycle maps commute
with the same pullbacks, proper pushes and determinant projectors.
At good finite-field fibres the degree-$2n$ Weil realization has
weight$2n$, and its Tate twist$(n)$ has weight zero, with the
purity locator already given in Section\ref{sec:framework}.
Thus(F1F1.2) and(F1F4.2) have the same integer factors in those
realizations. No estimate of a motivic Euler factor is used to
obtain degree8 or36.

In the project's archive notation, the actual bounded integral
representatives enter conditionA, cote16 scan050, through
SIGN\_NEXT SN-C03--C04 and the carried bounded-cycle reduction.
Their positive Weil Gram pairing is conditionI, cote16 scan049
and cote44 Theorem4. The integral graph polynomial $Re_W$ is
the Cayley--Hamilton correspondence inverse of cote16 scans048--049.
The carried HD-D03 extraction applied to degree8 or36 gives
$285120\,\delta\lceil\sqrt T\rceil$ or
$5773680\,\delta\lceil\sqrt T\rceil$, respectively, by
$4455\cdot8^2$ and $4455\cdot36^2$. The exact full class
from HD-F1-S07 permits the stronger extraction HD-F1-T05:
three actual unit pullbacks give multiplier$3\delta$ and
$c_t\le17\delta\lceil\sqrt T\rceil$ throughout that
principal component. Its integral lattice and effective-part
maps are displayed in Section\ref{sec:f1-exact-division}.

The additional original norm-image morphisms, its flat whole class,
Hilbert obstruction map and principal quotient are proved in
HD-F1-S09--S11 and T07. The terminal tensor exclusions T08--T09
retain the actual quotient curves, primitive projectors, paired
period maps and Grothendieck-group operations displayed there.
R09--R11 give the Fourier-kernel, Poincare, Adams and actual
truncation maps for the real-multiplication trace obstructions;
their complete common-annihilator computation takes place in
the whole Hochschild realization.


The further statements HD-F1-S12--S17, T10--T12 and R12--R15 retain
these exact realizations. The actual norm-fourfold Abel pushforward
has the whole GRR character of S14 and is proved to be the ordinary
structure sheaf of its normal image in S16--S17. Its formal neighbourhoods
use the actual scheme fibre ideals and the theorem on formal functions;
its unresolved trace is stated on this original object.
The terminal tensor closure T12 uses equality of the original rational
matrix-coefficient idempotents, finite quotient pullback/norm and explicit
Rosati matrix units; each polarization multiplier has its rank-six norm
written out. Thus its variational local-system maps and Poincare
biextensions carry the prescribed factor. R12--R15 use actual canonical
triangles, filtered diagonal-kernel HH evaluation, isogeny trace retractions
and universal Newton/Adams identities. Their exact framework locators are
F01--F04 and the source paragraphs of the respective proofs.


For S18--S19 the new maps are the derived quotient unit and multiplication,
the actual Picard dual isogeny, its eigensheaf algebra, and the finite
quotient with its resolution base-change square. For R16--R17 they are
the actual truncation arrows, the Yoneda extension-to-tangent map and
the differential of the coherent point action, compared by naturality
before the exact Euler pairing. For T13--T15 every map of the cover,
primitive Jacobian factor, restricted theta polarization, Poincare
biextension and variational realization is the identity; the subgroup
arguments bound ranks inside its original coefficient representation.
These are the exact additional realizations under the locators F01--F04.


For R18 the additional FGA maps are the invariant endpoint injections
of the original canonical triangle and its actual nonscalar endpoint
change map $d_\delta$. Its image is exactly the kernel of the written
extension-to-tangent map. The SGA~6 realization is additivity along
that triangle and the actual saturated stable-factor filtrations;
double dual and the specified big-open restriction/pushforward inverse
control their Hom spaces. The effective-cycle obstruction uses the
actual quotient sheaves $L/F_i$. These operations retain the exact
F01--F04 locators and introduce no change of object.


For T16 the only extra coefficient map is the actual surjection from
the simply connected finite group onto the original perfect central
layer; its character inflation and the displayed commutators bound the
original inertia ranks. For T17 the extra maps are the averaged Hom
projector, its isotypic evaluation and the centralizer quotient of the
original product-one tuple. All geometric, biextension, cup-product and
Adams maps remain identities. For S20 the displayed pluricanonical
pullbacks and Proj construction give an actual free action on $T$.
The object-supplied scalar cochain gives an explicit dg equivalence
of coherent actions; the ensuing torsor descent identifies the full
categorical traces. These are the additional exact realizations at
the F01--F04 locators, with no numerical replacement of an object.


\paragraph{Realization of the subsequent restrictions T18--T31, R19--R30 and S21--S30.}
These proved restrictions retain the exact locators F01--F04.
For T18--T24 and T26--T31 the finite central-layer maps,
intertwiners and polynomial projectors are the ones displayed
in their proofs, acting on the original coefficient local system.
The associated Jacobian corner, its restricted theta line bundle
and its cycle correspondences remain the original objects.
T25 additionally constructs its actual scalar-square quotient
and subgroup coefficient; its induction intertwiner, pairing
factor two and prescribed-corner polarization multiplier are
the displayed maps of HD-F1-T25 at their stated scope.
For the actual centralizer quotient in T31 the two Jacobian maps
are pullback and norm, whose composite is multiplication by
$|C_R|$; cup pairings acquire that same degree. Applying
(F1F3.1)--(F1F3.2) to this very quotient gives its Poincare
biextension and polarization comparison. The separate retained
sixfold corner has the identity comparison. Thus the finite-group
rank exclusions restrict the actual Hurwitz-to-period morphism
of F01 and the detector-source gate, while the quotient-cover
parameter count in T31 concerns precisely its displayed quotient
family.

For R19--R30 the objects are the original canonical triangle,
its endpoint sheaves, their stable filtrations, and the actual
hull and evaluation quotients. Their coherent $K_0$ equalities
are induced by those triangles and short exact sequences.
The degree-two leading Chern characters of the quotient sheaves
are their effective support cycles with generic-length
multiplicities. In particular the two cycles used in R29 obey
\[
 [A^{**}/A]_2-
 \bigl([I^{**}/I]_2+[Q^{**}/Q]_2\bigr)
                 =-\ell^2+q(A_1+A_2)
\]
in the three-row case, by its actual two quotient sequences.
Their supports define the corresponding Hilbert points and their
leading cycles define Chow points under F02. R29's period-lattice
Albanese homomorphism and R30's quotient $X\to X/B$ are actual
abelian-variety maps, so (F1F3.1) applies to them as well.
R30's affine restriction is followed by the proved inverse on
self-Ext through its closed-support comparison; the final Serre
transpose is the same original quotient map. These are the
specified FGA and duality realizations of its rank contradiction.

The tensor identities and field projectors in R26--R28 are
identities in the displayed rational cohomology of these actual
objects. Their underlying $K_0$ equations remain the preceding
triangle and quotient equations. Applying the Newton operation
$\psi^a$ to those equations is additive and multiplies every
degree-two leading cycle term by $a^2$ through (F1F4.1).
This is the precise SGA~6 graded realization. The Ext and
semiregularity restrictions remain attached to the original
objects and connecting maps, exactly as in (F1F4.5).

For S21--S29 the marked divisor family is the actual closed
Hilbert curve constructed in S26 and recovered in S28.
Its universal incidence subscheme determines the retained
symmetric divisors, including their multiplicities. The canonical
relative Proj of S25 and the actual Picard and Albanese maps
recover the original cyclic cover in S27--S29. Consequently
on the original geometric open the actual normal obstruction
has kernel equal to the image of the Prym derivative and image
of dimension $16-\operatorname{rank}\mu_C$. The whole S14
character is horizontal on the specified polarized Weil
tangent, by the determinant and polarization contraction
calculation. Under F01--F02 the quotient of that tangent by
the recovered-cover image is therefore sent by the stated
ideal boundary into the full trace kernel. These are actual
Hilbert and perfect-object obstruction maps for the same
horizontal de Rham character.

For S30 use its two actual finite isogenies, without changing
their domains or codomains. The torsor fibre-product splitting
and the unit/trace retractions give the stated self-Ext maps;
the formal first-neighborhood isomorphism gives their Atiyah
comparison. Their whole-character realizations are
\[
 \operatorname{ch}(f_*E)=f_*\operatorname{ch}(E),
       \qquad
 \operatorname{ch}(g^*E)=g^*\operatorname{ch}(E),
\]
with the trivial ambient Todd classes. The normalized Poincare
isomorphism (F1F3.1) applies to $f$ and $g$; the specified
polarizations satisfy $f^*\theta_A=\beta$ and
$\theta_B=g^*\beta$. The finite line-action algebra and
translation torsor provide the actual coherent group actions.
Their finite orbit sums have the diagonal obstruction maps
and signed $K_0$ classes written in S30. Applying (F1F4.1)
to each whole character gives its Adams realization degree
by degree, with the same finite pullback or proper pushforward.

Every actual finite collection of these algebraic objects,
endomorphisms, quotients, kernels and incidence subschemes has
the finite-presentation spreading map of EGA IV$_3$,8.8.2
already specified in F04. Their cycle and $K_0$ identities
therefore have the same Betti, de Rham and good-fibre
$\ell$-adic realization maps. The new exclusions use the
complex-fibre tangent, positivity and trace calculations
proved in their own sections. For the project's named rows,
these are explicit restrictions on DEG\_L Q-T7-REPLACEMENT
and the remaining source gates for Q-BOUNDS and Q-DIV;
they preserve the proved detector bounds and the original
condition A, condition I and Cayley--Hamilton maps at the
archive locators above.


\paragraph{Realization of the restrictions T32--T33 and S31--S33.}
The maps in T32 are the actual central-cover surjection
$\beta:H\to L$, its commuting square with each original
conjugation automorphism, and the displayed multiplicity-space
projectors. In T33 the scalar-character idempotents and the
finite Fourier operator act on the actual finite function space;
evaluation identifies the original coefficient space with
$W\otimes\operatorname{Hom}_H(W,V)$. Its formula
$\rho(g)=A_g\otimes T_g$ retains all scalar phases and
multiplicities. These comparisons estimate the original inertia
operators. The original curve cover, coefficient local system,
middle extension, cup form and prescribed Jacobian corner are
the same objects throughout. Their comparison morphisms are
identities, including the restricted theta homomorphism and
the normalized Poincare biextension in (F1F3.1).
Consequently their Hodge and de Rham realizations, period
derivative and good-fibre $\ell$-adic realizations retain those
identity maps. The bounds restrict the original source of the
Hurwitz-to-Weil period morphism in F01. In the named rows they
restrict the source gate for DEG\_L Q-BOUNDS and Q-DIV; no
additional cycle component is introduced by these comparisons.

For S31--S32 the actual object is
$F=E\otimes^{\mathbf L}N$ on the original abelian eightfold.
Its global coevaluation, the Koszul-signed evaluation, and
the first-jet Baer pushout supply the actual map
\[
 p_{\rm ns}\operatorname{ptr}_N
       \operatorname{ev}_F(\kappa+\eta)
       =-r\operatorname{ob}_{\rm jet,E}(\kappa),
 \qquad r=\operatorname{vrk}(N).
\]
The full-character identity and cup product with $\beta$ force
$\lambda t=rw\ne0$. These are maps between the deformation
and obstruction spaces of the displayed actual perfect objects.
The original Hilbert support and the ideal-boundary morphism
of S24 remain their F02 realization. For S31's concrete
isogeny-line objects $N=h_*L$, the proper pushforward character
is $r\exp(D)$ with
$D=(\deg h)^{-1}h_*c_1(L)$, through the proved pullback and
trace equalities. Their Poincare comparison is precisely
(F1F3.1) for this $h$. The ambient polarization and the
original norm support stay fixed.

S33 constructs the additional morphism for an arbitrary actual
coherent $\mathcal G$ on the normal integral norm image:
\[
 \begin{split}
 \mathcal T_{\mathcal G}:\quad
 \operatorname{ev}_{i_*\mathcal G}
     \bigl(D_P\oplus H^2(P,\mathcal O_P)\bigr)
       &\longrightarrow W_E,\\
 \operatorname{ev}_{i_*\mathcal G}(\kappa+\eta)
       &\longmapsto-r\operatorname{ob}_{\rm jet,E}(\kappa).
 \end{split}
\]
The preceding support-ideal construction proves its kernel
inclusion and surjectivity. In FGA terms its local ideals
$I_{\varepsilon,i}$ glue to an actual first-order Hilbert
point over the specified ambient deformation. Conjugation
removes the scalar triple-product defect before these ideals
are glued. Their normal-boundary class is exactly the one
mapped to $W_E$ in S24. Thus this is an explicit relation
between the scalar-twisted perfect-object deformation and
the ordinary Hilbert deformation of its normal support.
The full S14 character is horizontal under the stated HKR
shear, and the restriction of $-\mathcal T_{\mathcal G}/r$
maps its evaluated full-trace kernel onto $W_E$.
This is the F01--F02 obstruction statement for the same
horizontal de Rham character. It excludes these actual
candidates in DEG\_L Q-T7-REPLACEMENT at precisely the
scheme-theoretic quotient action $I_Y(i_*\mathcal G)=0$.

The exact coherent sequences and perfect triangles give the
corresponding identities in $K_0(P)$. For S32,
$[F]=[E]\,[N]$, with all differentials retained in the
actual object supplying that class. For S33 the generic-basis
sequence gives $\operatorname{ch}_{<4}(i_*\mathcal G)=0$
and $\operatorname{ch}_4(i_*\mathcal G)=r[Y]$ through
the localized support-character map. A zero-dimensional
ideal modification has its actual identity
\[
 [i_*\mathcal I_{Z_0/Y}]=[E]-[i_*\mathcal O_{Z_0}].
\]
Its top Chow term is $i_*[Z_0]$. The explicit cycle-class
map to $H^{16}(P,\mathbb Q)$ sends that term to
$\operatorname{length}(Z_0)\beta^8/\int_P\beta^8$.
Applying the Newton operation of SGA~6 to these actual
$K_0$ identities gives (F1F4.1): in particular the
codimension-four term is multiplied by $a^4$ and the
point term by $a^8$. The self-Ext and characteristic-image
maps remain the displayed maps of the actual original
objects; no object-level semiregularity conclusion is
deduced from an equality of $K_0$ classes.

The exact original F01--F04 locators, including the stated
SGA~7 VIII scan limitation, remain unchanged. EGA
IV$_3$,8.8.2 applies to the finite presentation data of
these actual coherent objects, maps and ideal sequences
as specified in F04. Their cycle and $K_0$ identities
therefore have the same Betti, de Rham and good-fibre
$\ell$-adic realization maps. Their complex tangent and
trace exclusions use the proofs in S31--S33 and T32--T33
at the stated scopes. The previously proved condition A,
condition I, Cayley--Hamilton projector and integral
division bounds are retained at their archive locators.


\paragraph{Realization of T34--T36, the two individual Hurwitz strata, and S34--S35.}
The representation maps of T34 are the actual scalar-sector
idempotents, the specified field-semilinear descent with its
Hermitian multiplier, and the original coefficient commutator.
T35 identifies the literal Clifford involution and constructs
the descended map $\overline\beta:\Omega\to L$ whenever
that involution is killed. The automorphism squares commute
with the same quotient. T36 then uses the actual isotypic
projections and the evaluation isomorphisms
\[
 V_{\mathbf C}=\bigoplus_i W_i\otimes
                 \operatorname{Hom}_H(W_i,V_{\mathbf C}).
\]
Its displayed injections into bounded tensor powers carry
the original $G$-action to their image by
$\iota\rho(g)\iota^{-1}$. Each branch estimate concerns
the original inertia operator on $V$. The original cover,
local system, middle-extension cohomology, cup product,
primitive Jacobian image and restricted polarization have
identity comparisons. Their pairing multiplier is one.
The normalized Poincare comparison (F1F3.1), Hodge
filtration and period derivative therefore also retain
the same identity maps. These restrictions narrow the
actual Hurwitz source in F01 and introduce no additional
cycle component into the condition A or division bounds.

For HD-F1-I01--I02 the new morphism is evaluation of
the actual $\mathcal O_Y(Q)$ on the scheme-theoretic
fibres of the same degree-nine map $h:Y\to\mathbf P^1$.
In the four-value case this identifies
$h_*\mathcal O_Y(Q)$ with $\mathcal O^9$ and carries
the literal unit to a constant basis vector. In the
one-reflection case the exact kernel sequence is
\[
 0\longrightarrow\mathcal O(-1)
   \xrightarrow{(X,Y,0,\ldots,0)^t}\mathcal O^{10}
   \longrightarrow h_*\mathcal O_Y(Q)\longrightarrow0.
\]
The written quotient map $(-Y,X)$ and triangular
bundle automorphism retain the specified unit and give
$\mathcal W_Q=\mathcal O(1)\oplus\mathcal O^7$.
The I21 pushout and cohomology map identify the actual
forgetful tangent kernels with
$H^0(\mathcal W_Q(2-\nu))$; the two calculations give
period ranks one and two respectively. These are
relative coherent evaluation and restriction maps,
with the EGA III$_2$,7.7.5--7.8.4 base-change framework
at the previously checked scope. The polynomial
multiplication of the two primitive eigendifferential
spaces and the canonical quadratic trace are the
actual period codifferential.

The comparison with the original Jacobian corner is
the same restricted pullback $q^*$ and norm $q_*$,
with $q_*q^*=[|\ker\chi|]$. The restricted theta
pullback has precisely that multiplier. Formula
(F1F3.1) applied to this actual homomorphism is its
Poincare realization; its rational symplectic
similitude gives the same period-rank comparison.
The two full-group tuples give actual connected
covers through the retained finite-cover
correspondence. Their marked Hurwitz maps also
give the precise comparison with Filom's published
nonvanishing theorem recorded in the proofs.
The evaluated ranks concern these individual
strata and add no effective detector cycle.

For S34 the object is the actual
$F=Ri_*\mathcal G$, with every differential and
coupling in $\mathcal G\in D^b_{\rm coh}(Y)$ retained.
Its epsilon triangle and derived-adjunction
factorization $\mu_f=j b_f p$ yield the intrinsic
regular alternating trace. The proved trace-ambiguity
calculation gives the actual support ideals whose
normal cocycle is the original S24 cocycle.
Thus its characteristic-image quotient is
\[
 \mathcal T_{\mathcal G}\bigl(\operatorname{ev}_F
            (\kappa+\eta)\bigr)
                =-r\operatorname{ob}_E(\kappa),
 \qquad r=\sum_a(-1)^a\operatorname{rk}_Y H^a(\mathcal G).
\]
The complete character forces $r\ne0$, with either
sign allowed. For S35 the actual sheaf is
$F=i_{k*}V$ on $Y_k=(Y,\mathcal O_P/\mathcal I^k)$.
Its ordinary locally free matrix trace constructs
the thickened support ideals. The actual sheaf map
\[
 \rho_k:\mathcal Hom_P(\mathcal I^k,\mathcal O_P/\mathcal I^k)
       \longrightarrow
 i_*\mathcal Hom_Y(\mathcal I/\mathcal I^2,\mathcal O_Y),
 \qquad \rho_k\partial_k=\partial_1
\]
maps their normal cocycle to the reduced one.
Its explicit normalization is
$r_k=\binom{k+3}{4}$. Consequently
$\mathcal T_{k,V}(\operatorname{ev}_F(\kappa+\eta))
=-m r_k\operatorname{ob}_E(\kappa)$.
These are actual relations between perfect-object
deformation and Hilbert ideal deformation in F02.
Their evaluated full-character shear surjects onto
the same $W_E$ and is fixed by every specified
coherent finite point lift. This gives the proved
trace exclusions for DEG\_L Q-T7-REPLACEMENT.

The exact truncation triangles give
$[Ri_*\mathcal G]=\sum_a(-1)^a[i_*H^a(\mathcal G)]$
in $K_0(P)$ without splitting the original object.
The localized leading-cycle and cycle-class maps
give $\operatorname{ch}_4(F)=r[Y]$ for S34 and
$\operatorname{ch}_4(F)=m r_k[Y]$ for S35, with
lower components zero. In the second identity
$m r_k$ is the actual generic length of the
nonreduced module. Applying the Newton-defined
Adams operation multiplies this codimension-four
term by $a^4$, as in (F1F4.1). The characteristic
and self-Ext maps remain those of the original
objects in the displayed proofs.

\paragraph{Realization of the pure RM cancellation restriction R31.}
The new geometric maps are the resolution of each
actual surface component, its Albanese homomorphism,
the two actual curve fibrations in the excluded
four-dimensional-image case, the induced Jacobian
product map, and the isogeny inverse obtained by
descending multiplication. The resulting equality
$\gamma a=[M]\psi$ proves the required relation
on the same holomorphic one-forms and rational
homology. The projective quotient by each surviving
abelian surface is the actual quotient in R30;
its lattice and polarization maps retain the
Poincare formula (F1F3.1).

On the coherent side R31 retains the original
$p:A\to Q$, $l:Q\to L$ and connecting morphism
$\delta:T\to A[2]$. Its actual exact sequences give
\[
 \begin{aligned}
 [E]&=[A]-[T],& [Z]&=2[L]-[T],\\
 [J]&=[L]-[Q],& [Z]&=[Z_0]+[P],\\
 [P]&=\sum_a[P_a]-[C_0].
 \end{aligned}
\]
The last equality is used in the excluded
zero-pure-coefficient case and keeps every
component thickening. The lower-dimensional
terms disappear from the degree-two Ext comparison
through both proved adjacent vanishings.
The support cycle is the leading localized
character of the actual $Z$, followed by its
cycle-class map. The two positive pure coefficients
are computed by pairing that cycle with the
opposite semipositive RM block form.

The original quotient columns, their null
cohomology action after affine base restriction,
and their actual Serre transpose identify the
precise predecessor image. Full Atiyah-trace
cyclicity gives the same original endpoint
Yoneda coefficient. The block character identity
then proves
\[
 \beta_L=l[2]p[2]\delta=0,
 \qquad
 0\ne\beta_Q=p[2]\delta=\partial\xi,
 \quad \xi\in\operatorname{Ext}^1_G(T,L/Q).
\]
The four-dimensional pure RM subspace belongs
to the actual invariant predecessor image.
The source correction from
$\operatorname{Ext}^1_G(\ker p,Q)$ remains in
the same Hom sequence. These are F02 perfect-object
and F04 coherent $K_0$ maps for the stated
replacement gate. The Adams action multiplies
their codimension-two support character by $a^2$
and their divisor character by $a$; no statement
about the remaining Ext groups follows solely
from these graded character identities.

All exact F01--F04 source locators and their
recorded reading limits remain in force,
including the original SGA 7 VIII scan limitation.
EGA IV$_3$,8.8.2 applies to the finite lists of
actual algebraic maps and coherent data, with
the Betti, de Rham and good-fibre etale
realizations already specified. The tangent
and trace assertions here retain their proved
complex-characteristic scope. No new analytic
Euler-factor estimate, effective-cycle bound
or integral division bound is inferred from
this realization record.


\paragraph{Realization of the coherent thickening exclusions S36--S37.}
For S36 the object is the actual coherent module on the ordinary
power$Y_k$, including its smaller-dimensional torsion and defects.
Its generic matrix trace is followed by the actual normal retraction
$\rho_k$. The finite resolution over the completed support coefficient
ring with dual numbers makes the whole multiplication trace regular.
Faithful flat contraction and the normal-domain intersection descend
that same rational normal value to$Y$. Thus the displayed local
ideals give actual flat Hilbert deformations of the original reduced
support. The characteristic-image quotient is
\[
 \mathcal T_F\bigl(\operatorname{ev}_F(\kappa+\eta)\bigr)
       =-m\binom{k+3}{4}\operatorname{ob}_E(\kappa).
\]
S37 uses the actual generic module isomorphism
\[
 F_\xi\simeq\bigoplus_a(\mathcal O_\xi/\mathcal I_\xi^a)^{m_a}.
\]
Its new map takes the epsilon coefficient of multiplication by
$f\in\mathcal I^a$ on a lifted central vector and applies the
actual central projection$\pi_a$. The resulting map$b_f^a$
is independent of that lift because$\pi_a(fw)=f\pi_a(w)=0$.
The diagonal block of the commuting-matrix tangent equation
constructs its exact local moment trace. Every off-diagonal
first-order matrix is retained; a return across blocks requires
two epsilon factors. The total trace therefore gives the
decomposition-independent reduced normal functional, with
\[
 \ell=\sum_a m_a\binom{a+3}{4},\qquad
 \mathcal T_F\operatorname{ev}_F(\kappa+\eta)
                       =-\ell\operatorname{ob}_E(\kappa).
\]
The same actual finite resolution proves regularity at the defects.
Transport of the central projections through the original gluing
maps gives the normal cocycle$\partial_1\kappa$; scalar triple
products act trivially by conjugation. This is the precise
F02 map from the original perfect-object deformation to the
Hilbert ideal deformation, followed by its original obstruction.

The localized coherent cycle and cycle-class maps give
$\operatorname{ch}_4(F)=\ell[Y]$ and zero lower components.
The complete S14 equality retains all higher components and
identifies$\lambda t=\ell w$. Its full sheared characteristic
image surjects onto$W_E$, with each specified coherent point
lift retained. For the actual two-point kernel of S37, coherent
$K_0$ gives$[M]=[\mathcal O_{Y_2}]+[\mathcal O_Y]-[\mathbf C_p]
-[\mathbf C_q]$. The point terms become the stated top
polarization class only after cycle class. Newton-defined
Adams operations multiply the leading codimension-four term
by$a^4$ and the point terms by$a^8$, through(F1F4.1).
These are the actual F04 maps behind the trace exclusions.

\paragraph{Realization of the two further nine-letter period ranks I03.}
For the original finite map$h$ and its reduced uncolored divisor$Q$,
the evaluated positive subbundle and its actual quotient give
\[
 0\longrightarrow H^0(h_*\mathcal O(Q)(-1))\otimes\mathcal O(1)
 \longrightarrow h_*\mathcal O(Q)
 \longrightarrow\mathcal O^{9-A}\longrightarrow0,
 \qquad A=2,3.
\]
The unit is the prescribed section with zero divisor$Q$ under
the actual polynomial presentation. Its nonzero constant quotient
component, proved from the homogeneous identity for$h$, permits
the stated constant and triangular automorphisms. Thus the original
trace quotient is$\mathcal W_Q=\mathcal O(1)^A\oplus
\mathcal O^{8-A}$. The I21 pushout and cohomology maps give
$H^0(\mathcal W_Q(-2))=0$, and the actual period ranks are
$A+1$. Evaluation is on each scheme-theoretic fibre, including
ramification. The two full-group tuples give actual connected
covers in the retained finite-cover correspondence.

The actual polynomial product of primitive eigendifferentials
and quadratic trace is the period codifferential. The restricted
$q^*$ and$q_*$ satisfy$q_*q^*=[|\ker\chi|]$, and the theta
pullback has exactly that multiplier. Applying(F1F3.1) to this
same$q$ gives its Poincare biextension comparison. Consequently
the original local system, Hodge filtration and period derivative
have the stated equal rank. This is the exact F01 variational
realization of the two component computations.

\paragraph{Realization of the odd symplectic restrictions T37--T38.}
The universal-cover map$\beta:H\to L$ and the lifted outer map
obey$\beta\alpha_g=c_g\beta$. T37 constructs actual Weyl-algebra
isomorphisms, their genuine unitary lifts and parity projectors.
Its scalar correction on the natural finite-field space induces
the same$\alpha_g$. On each original isotypic coefficient space,
evaluation identifies the original operator with the specified
intertwiner tensored with its actual multiplicity operator;
every coefficient-space phase stays in that latter operator.

T38 instead maps an actual transvection to the original
commutator$\beta(\alpha_g(t)t^{-1})=[g,\beta(t)]$.
The delta-function tensor map and the normalized coset-sum
embeddings into$\mathbf C[\operatorname{Hom}(X,U)]$ provide
the exact auxiliary representation comparisons. The full
orthogonal-group central projector is divided by its actual
multiplicity in the trace formula. The stable lower-rank code
maps and top-weight mutual commutants then bound the trace
on each original irreducible constituent. The original
isotypic evaluation sums these bounds with all multiplicities.
This proves the stated finite natural-rank cutoff
$(3\cdot10^{30}-3)^2+4$ for this branch.

These comparisons estimate matrices in the unchanged original
coefficient system. Its cover, parabolic cochains, primitive
Jacobian corner, integral lattice and cup polarization have
identity comparisons, with pairing multiplier one. Their
de Rham, Betti and good-fibre etale realizations therefore
retain the same graph correspondences and Poincare map.
They restrict the F01 Hurwitz source and add no effective
detector cycle or new division constant.

\paragraph{Realization of the actual RM triangle exclusions R32--R34.}
R32 retains the two original quotient columns and their
injective hull map. A square-zero endomorphism factors through
its actual coherent image; the reverse composite is zero.
Atiyah naturality, cyclic trace and the full Hochschild module
identity therefore give the actual vanishing used to prove
ordinary simplicity. This is an object-level trace statement
on the existing quotient target.

R33--R34 retain the same actual triangle and its coherent
$K_0$ relation$[E]=[A]-[T]$. Their full exterior maps
$\Phi_i$, the invertible$L_{\xi_i}$ and the operators
$P_i=L_{\xi_i}^{-1}L_{\alpha_i}$ act on its actual
character-map spaces. R34's quadratic operator identity,
eigenprojectors and repeated-root image/preimage bases
give the complete tensor nullity estimate. Naturality
for the given point lifts factors each full character
map through the original invariant endpoint Ext group.
Their rank sum is at least$24$ in R33 and$28$ in R34.

The actual Yoneda maps$\nu(f)=\delta f$ and
$\mu(f)=f[2]\delta$ use the original coupling throughout.
Their triple-composition map$J_\delta$ is paired by the
actual invariant Serre trace. The twice-rotated TR3 diagram
maps the space of triangle morphisms to its exact endpoint
compatibility kernel. Its middle projection has dimension
loss$2$ in R33 and at most$4$ in R34; the endpoint
compatibility loss is respectively$2$ and at most$3$.
The resulting contradictions with the original full balanced
rank$18$ exclude the complete rank-$(3,2)$ equal-slope
semistable TF adjacent gate. Every coherent surface and
lower-dimensional defect is retained in those maps.

These are precise F02 perfect-object restrictions for
DEG\_L Q-T7-REPLACEMENT. Their leading coherent surface
cycles and the original determinant tensor remain the F04
cycle-class and$K_0$ maps. Adams operations act on their
graded characters by(F1F4.1); the Ext contradictions
belong to the actual original triangle, whose entire
morphism spaces have been kept. The broader RM seed
problem remains outside the excluded ranks and hypotheses.

The exact Grothendieck locators in F01--F04 remain those
already audited: the 1966 de Rham formulation and comparison,
FGA 221 Theorems 3.1--3.2 and Proposition 3.8, SGA 7 VII 2.9.5--2.9.6
and VIII 3.2, and SGA 6 V Appendix 7.1--7.5. The recorded
original VIII scan limitation is unchanged. EGA IV$_3$,8.8.2
spreads the finite lists of actual algebraic maps and
coherent data; the tangent and trace conclusions retain
their proved characteristic-zero scope. In the archive
framework, the same cycle/projector maps feed conditionA
at cote 16 scan 050, condition I at scan 049 and cote 44
Theorem 4, and the Cayley--Hamilton inverse at scans 048--049.
This realization record introduces no additional strict
movement, effective-cycle bound or analytic L-value estimate.


\paragraph{Realization of the unrestricted low-Hom RM restriction R35.}
The object remains the given equivariant triangle
$A\to E\to T[-1]\xrightarrow{\delta[-1]}A[1]$.
Its actual$K_0$ identity gives$[E]=[A]-[T]$ and, through
the same rational cycle-class map, the entire balanced character
difference. The source$B$-shear and target multiplication by
$e^B$ are the displayed invertible cohomology maps. They
preserve the original full character-map ranks and never
postulate an object with a chosen formal tensor.

The two projections from the actual triangle-morphism space
give the precise inequality
\[
 e^2_G(E,E)\ge e^2_G(A,A)+e^2_G(T,T)
       -r_D-r_\nu-r_\mu+r_J,
 \qquad r_D\le h,\quad r_\nu,r_\mu\le h.
\]
The complete restricted-character classification applies at
all positive coherent endpoint ranks. The actual product
pairing gives$r_J\ge2$. The complex pencil projectors and
Jordan bases then leave exactly the three rows in R35.
The two zero-coefficient rows imply the full character
$R\exp(\ell-B)$ of an existing endpoint, through its
degree-two component and the exact Mukai pairing identity.
These are F02 restrictions on the original perfect-object
deformation and F04 restrictions on its actual character
realization. The associated Hilbert/Chow classes, their
Newton-defined Adams operations and the original cycle maps
retain the F01--F04 domains and normalizations.

\paragraph{Realization of the supplementary explicit tuples I04.}
Every listed branch tuple is a literal product-one generating
tuple in the indicated monomial matrices. Its finite-cover
realization and its subgroup quotient give the original
map$h:Y\to\mathbf P^1$ and the cyclic cover above it.
Every merging or splitting is the displayed adjacent matrix
product, and preserves the calculated quotient genus and
both primitive eigenspace multiplicities. The exact identity
$N+\nu=2m+2n-A$ identifies its rank range with the already
proved I16 range. The I16 jet map, I21 unit pushout and
I20 canonical quadratic trace are the same variational maps,
and the actual$q^*$ and$q_*$ give the F03 comparison with
their specified theta multiplier. The finite covers and
their marked-divisor Hilbert realization therefore give
nonempty components with the stated period ranks.
The additional assertion is uniform tuple existence;
the inherited rank criterion is not counted again.

\paragraph{Realization of the selected-cohomology restriction S38.}
The actual scalar-twisted perfect lift retains all differentials
and all extension maps. Its epsilon triangle induces the two
ordinary Bockstein maps adjoining the selected cohomology$M$.
Their vanishing gives the actual flat sheaf$\widetilde M$,
whose ideal-action trace in S37 produces the flat embedded
deformation of the reduced$Y$. The resulting quotient is
\[
 \mathcal T_F\operatorname{ev}_F(\kappa+\eta)
                  =-\ell_M\operatorname{ob}_E(\kappa),
 \qquad\ell_M=\sum_a m_a\binom{a+3}{4}>0.
\]
This is the explicit F02 map from the original object
deformation to the Hilbert ideal deformation. It yields
the full invariant trace-kernel bound on the stated S14
class. For the actual three-point-ideal example, the
coherent$K_0$ relation is
$[F]=[E]-[k_p]-[k_r]+[k_q]$.
Only after rational cycle class do these point terms
become the top polarization class. The support-Ext image
is compared by$\Psi_{q,s}$ and the scalar unit retraction;
the chosen coupling has nonzero reduced normal image and
zero scalar retraction. Those maps prove the failure of
a support presentation for the actual indecomposable cone.
Adams operations act on these same classes by(F1F4.1).

\paragraph{Realization of the even symplectic restriction T39.}
The natural central cover and each lifted automorphism satisfy
$\beta\alpha_g=c_g\beta$. The displayed transvection
commutator maps under$\beta$ to the original coefficient
commutator. The actual unitary-Heisenberg coordinate
homomorphism, the genuine quadratic models, the explicit
norm-plane isometry and every scalar Fourier phase identify
the original restricted constituent with its actual tensor
summand. The normalized coset embeddings and the explicitly
rescaled coset subtraction preserve their exact maps.
The full orthogonal-group projector divides its trace by
the actual multiplicity, and its evaluated multiplicity
space injects into the original quadratic module.
These give the proved trace bound for that original
commutator and the unchanged inertia-rank contradiction.

On the original cover, parabolic cohomology, primitive
sixfold corner and Hermitian cup pairing, the comparison
is the identity with polarization multiplier one.
Thus the F01 variational source, F02 Hilbert/cycle map,
F03 Poincare biextension and F04 Adams/cycle realization
are carried through their original maps. The result
restricts possible source tuples to the stated finite
natural-rank range and supplies no new detector cycle.


\paragraph{The codimension-two cohomology extension S39.}
The original Bocksteins vanish on$P\setminus Z$ through
their stated ordinary restriction maps. The selected
cohomology is flat on that open. The actual normal trace
there gives conormal values in$\mathcal O_{Y\setminus Z}$.
Each finite ideal-generator value extends uniquely by
the normal-domain intersection over height-one local rings.
Every relation and ambient transition extends by the same
map. The explicit ideals$f+\epsilon h$, with
$h\bmod I=-u(f)$, are flat and glue in the original
ambient deformation. This is the F02 Hilbert comparison
also at the omitted set, retaining all complex couplings.
The new actual example has$\operatorname{Hom}(A,T)=\mathbf C^5$
supported at$p$, while its full character is obtained
from$[F]=[E]-2[k_p]-[k_r]+[k_q]$ by the same F04
cycle map. Its nonzero normal coupling and explicit
nilpotent endomorphism ideal prove the asserted scope.


The audited exact locators remain F01--F04: Grothendieck's
1966 de Rham formulation; FGA221 Theorems3.1--3.2 and
Proposition3.8; SGA7VII2.9.5--2.9.6 and VIII3.2; and
SGA6V Appendix7.1--7.5. The recorded VIII scan limitation
is unchanged. EGAIV$_3$,8.8.2 spreads the finite coherent
and algebraic data through these same maps; the proved
tangent and trace conclusions retain their characteristic-zero
scope. The original cycle/projector comparisons retain
conditionA at cote16 scan050, conditionI at scan049 and
cote44 Theorem4, and the Cayley--Hamilton inverse at
scans048--049. This realization record adds no separate
strict movement or analytic L-value estimate.


\paragraph{The local derived-annihilator comparison S40.}
The original bounded coherent complex and its actual perfect
first-order lift retain the epsilon triangle with maps$j,p$,
where$pj=0$ and$jp=\epsilon$. Off the stated codimension-two
set, multiplication by$f\in I_Y$ is zero on the original
object in the local derived category. The actual adjunction
therefore factors its action on the lift as$j b_f p$.
An ambiguity$j b p=0$ has zero intrinsic generic supertrace:
the displayed finite separable coefficient-field extension,
restriction to a perfect complex over$K[\epsilon]$ and
chain-homotopy trace give
$\operatorname{Tr}(jBp)=\epsilon\operatorname{str}(B)=0$.
Thus$r^{-1}\operatorname{str}(b_f)$, for the nonzero signed
generic cohomology rank$r$, defines the actual normal value.
The identity$pj=0$ kills products from$I_Y^2$.

The height-one regularity and the normal-domain intersection
extend these values, all conormal relations and all ambient
transitions across the omitted set. The explicit ideals
$f+\epsilon h$, with$h\bmod I_Y=-r^{-1}\operatorname{str}(b_f)$,
give the original flat embedded deformation. This is the
F02 map from the actual complex deformation to the Hilbert
ideal deformation, and its obstruction quotient is exactly
$-r\operatorname{ob}_E$. The leading-cycle equality
$\operatorname{ch}_4(F)=r[Y]$, followed by the same full
S14 character and point-action comparison, gives the
stated invariant trace-kernel bound. No global support
presentation is used in defining these maps.

The dual-cone example uses the actual perfect duality
$D=R\mathcal Hom_P(-,\mathcal O_P)$ and its bidual map.
The literal closed-immersion Hom adjunction gives, in the stated scope,
\[
 D(Ri_*H)\simeq Ri_*R\mathcal Hom_Y(H,\omega_Y[-4]).
\]
It transports a putative support presentation back to the
original cone, where its nonzero normal coupling is
detected by the reviewed support-Ext map and scalar
retraction. The example's full character follows from
these exact triangles and dual maps, before rational
cycle class is applied. F04 retains these original
$K_0$ classes, Adams operations and cycle maps.

\paragraph{The split orthogonal comparison T40.}
The actual natural cover and all lifted semilinear
similitudes satisfy$\beta\alpha_g=c_g\beta$.
The displayed root commutator maps to the original
coefficient commutator. Averaged Hom evaluation turns
the genuine character equality into an actual module
isomorphism; a nonzero coordinate projection embeds
the original constituent into its permutation tensor.
The partial Fourier map, paired root-group Weyl words,
normalized isotropic-code embeddings and their
rescaled coefficient subtraction are the explicit
intertwiners throughout the lower-weight argument.
The full-weight torsor supplies both regular actions,
and the evaluated multiplicity space gives the exact
projector with coefficient$1/|\operatorname{Sp}(U)|$.
All scalar phases and original restriction multiplicities
are retained in the original inertia-rank contradiction.

These maps restrict possible original source parameters.
The original curve and local system, parabolic cochains,
primitive Jacobian corner, Hermitian cup and restricted
polarization are carried by the identity, with multiplier
one. Their F01 period image, F02 cycle realization,
F03 Poincare biextension and F04 Adams/cycle maps
therefore remain those of the same source.

\paragraph{The prescribed-unit comparison I05.}
For the actual rational degree-nine map$h$ and reduced
uncolored divisor$Q$, finite-flat fibre evaluation gives
the original positive subbundle in$h_*\mathcal O(Q)$.
The prescribed unit is followed through the displayed
constant and triangular automorphisms and the map
$(u,v)\mapsto x_0v-x_1u$. These give the actual two
quotient-bundle sequences. In homogeneous coordinates,
their section map is
\[
 (\alpha,\beta)\longmapsto\alpha p_0+\beta p_1,
 \qquad q_Q\text{ is the prescribed unit column}.
\]
The I21 tangent-kernel isomorphism and I20 canonical
quadratic trace, composed with the onto multiplication
$H^0(\mathcal O(4))^{\otimes2}\to H^0(\mathcal O(8))$,
give the exact original period rank selected by that
augmented matrix. The F03 comparison is the retained
actual$q^*$ with$q_*q^*=[|\ker\chi|]$ and its theta
multiplier. These are the F01 and F02 maps on the
original Hurwitz and marked-divisor Hilbert spaces.

The four product-one tuples recover the original
monomial generators by literal products. The inner
diagonal matrix$J$ displayed in the review identifies
the independent checker coordinates with these same
tuples. Their finite-cover realizations have the proved
constant image ranks$5,6,7,8$. The lower bound$A$
is credited to Filom through the exact marked-source
quadratic trace. The additional polynomial selector
and four evaluated profiles retain the original vector
bundle$K_0$ identities and their F04 Adams/cycle maps.


\paragraph{The anisotropic coefficient comparison T41.}
The explicit line-form similitude and norm-plane
isometry conjugate the actual natural group and
transport$\beta$ by their inverse conjugations.
The norm-plane diagonal and Frobenius matrices retain
every original lifted automorphism. The genuine
smaller symplectic action is obtained by restriction
from rank$k+5$, followed by the actual inclusion and
evaluation on the last five coordinates. Its exact
Heisenberg covariance makes the large unipotent
operators act on the anisotropic coefficient space.
Their fixed vectors give the additional coefficient
condition in every normalized code embedding; the
zero-column inclusion identifies the smaller original
permutation tensor. On the full-weight fibre the
unitary map$e_s\otimes v\mapsto e_s\otimes\rho_0(s)^{-1}v$
exhibits both the complete coefficient matrix algebra
and the right regular stabilizer action. It proves
the top multiplicity and the precise original trace
projector. The resulting bound preserves the same
F01--F04 source maps and multiplier one as T40.
The odd-dimensional characteristic-two case retains
the explicit T36 polar-radical group isomorphism
and the T39 bound through that map.


The foundational F01--F04 locators and their recorded
reading limits remain those audited above. EGAIV$_3$,
8.8.2 spreads the finite algebraic data through the
displayed maps; all proved tangent and trace conclusions
retain their characteristic-zero scope. This realization
record adds no separate strict movement or analytic
$L$-value assertion.
